% ============================================================================
% 第三章：实验与结果分析
% ============================================================================

\section{实验设计}

\begin{frame}{实验环境与数据集}
    \begin{block}{数据集：BCI Competition IV 2a}
        \begin{itemize}
            \item 9 名健康被试，4 类运动想象任务 (左手、右手、脚、舌头)
            \item 22 通道 EEG，采样率 250 Hz
            \item 训练集/测试集各 288 个试次
        \end{itemize}
    \end{block}
    
    \vspace{0.3cm}
    \begin{block}{预处理流程}
        \begin{enumerate}
            \item 带通滤波 (8--30 Hz): 保留$\mu$和$\beta$节律
            \item 陷波滤波 (50 Hz): 消除工频干扰
            \item 平均参考：提高空间分辨率
            \item 时间窗截取：根据算法选择截取
        \end{enumerate}
    \end{block}
\end{frame}

\begin{frame}{特征提取与分类}
    \begin{block}{CSP 特征提取}
        \[
        J(\mathbf{w}) = \frac{\mathbf{w}^{\text{T}} \Sigma_1 \mathbf{w}}{\mathbf{w}^{\text{T}} \Sigma_2 \mathbf{w}}
        \]
        \begin{itemize}
            \item 4 个 CSP 分量
            \item 对数方差特征
        \end{itemize}
    \end{block}
    
    \vspace{0.3cm}
    \begin{block}{分类器}
        \begin{itemize}
            \item 线性 SVM ($C = 1.0$)
            \item Z-score 标准化
            \item 5 折交叉验证
        \end{itemize}
    \end{block}
\end{frame}

\begin{frame}{对比方法}
    \begin{table}[h]
        \centering
        \footnotesize
        \begin{tabular}{l|c}
            \hline
            \textbf{方法类别} & \textbf{具体方法} \\
            \hline
            固定时间窗 & FL-2s, FL-1.5s, FL-1s \\
            搜索优化 & Grid-Search, Random-Search \\
            强化学习 & Standard Q-Learning, Double Q-Learning, \\
            & Dueling Q-Learning, Dueling Double Q-Learning, DQN \\
            \hline
        \end{tabular}
    \end{table}
\end{frame}

\begin{frame}{评估协议}
    \begin{itemize}
        \item \textbf{重复实验}: 每名被试 5 次重复实验，报告平均值 $\pm$ 标准差
        \item \textbf{评价指标}:
        \begin{itemize}
            \item 准确率 (Accuracy)
            \item Cohen's Kappa 系数
            \item F1 分数
            \item 标准差 (Std Dev)
            \item 变异系数 (CV)
        \end{itemize}
        \item \textbf{统计检验}: 配对样本 $t$ 检验 ($\alpha = 0.05$) + Cohen's $d$ 效应量
    \end{itemize}
\end{frame}

% ============================================================================
% 基线方法结果
% ============================================================================
\section{基线方法性能分析}

\begin{frame}{固定时间窗 vs 搜索优化}
    \begin{table}[h]
        \centering
        \footnotesize
        \begin{tabular}{lccccc}
            \hline
            \textbf{Method} & \textbf{Acc (\%)} & \textbf{Kappa} & \textbf{F1} & \textbf{Std} & \textbf{CV (\%)} \\
            \hline
            Fixed-2s & 45.08 & 0.266 & 0.426 & 8.95 & 19.85 \\
            Fixed-1.5s & 45.05 & 0.266 & 0.425 & 10.19 & 22.61 \\
            Fixed-1s & 45.47 & 0.272 & 0.434 & 10.87 & 23.91 \\
            Random-Search & 61.95 & 0.492 & 0.611 & 8.38 & 13.52 \\
            Grid-Search & \textbf{63.35} & \textbf{0.511} & \textbf{0.625} & 8.78 & 13.86 \\
            \hline
        \end{tabular}
    \end{table}
    
    \vspace{0.3cm}
    \begin{alertblock}{主要发现}
        \begin{itemize}
            \item 搜索方法显著优于固定窗口 (63.3\% vs 45.0\%, +40.6\%)
            \item 变异系数降低约 30\% (13.5\% vs 19.8\%)
            \item Grid-Search 略优于 Random-Search (+1.4\%)
        \end{itemize}
    \end{alertblock}
\end{frame}

% ============================================================================
% Q-Learning 与基线对比
% ============================================================================
\section{Q-Learning 与基线方法对比}

\begin{frame}{Q-Learning 方法族性能对比}
    \begin{table}[h]
        \centering
        \footnotesize
        \begin{tabular}{lccccc}
            \hline
            \textbf{Method} & \textbf{Acc (\%)} & \textbf{Kappa} & \textbf{F1} & \textbf{Std} \\
            \hline
            Grid-Search & \textbf{63.35} & \textbf{0.511} & \textbf{0.625} & 8.78 \\
            Standard Q-Learning & 62.68 & 0.502 & 0.620 & \textbf{8.20} \\
            Dueling Q-Learning & 62.74 & 0.503 & 0.620 & 8.26 \\
            DQN & 62.70 & 0.502 & 0.620 & 9.21 \\
            Random-Search & 61.95 & 0.492 & 0.611 & 8.38 \\
            \hline
        \end{tabular}
    \end{table}
    
    \vspace{0.3cm}
    \begin{exampleblock}{结论}
        \begin{itemize}
            \item Q-Learning 方法与 Grid-Search 差距仅 0.6\%--1.6\%
            \item Cohen's $d < 0.2$，差异不显著
            \item Standard Q-Learning 稳定性最优 (Std = 8.20\%)
        \end{itemize}
    \end{exampleblock}
\end{frame}

% ============================================================================
% 表格型 Q-Learning vs DQN
% ============================================================================
\section{表格型 Q-Learning vs DQN}

\begin{frame}{整体性能对比}
    \begin{table}[h]
        \centering
        \begin{tabular}{lcccccc}
            \hline
            \textbf{Method} & \textbf{Acc (\%)} & \textbf{Kappa} & \textbf{F1} & \textbf{Std} & \textbf{Window (s)} \\
            \hline
            Standard Q-Learning & \textbf{62.68} & 0.502 & 0.620 & \textbf{8.20} & 0.88 \\
            DQN & \textbf{62.70} & 0.502 & 0.620 & 9.21 & 1.05 \\
            \hline
            $\Delta$ & \textbf{-0.02} & -0.0003 & +0.0001 & \textbf{-1.01} & \textbf{-0.17} \\
            \hline
        \end{tabular}
    \end{table}
    
    \vspace{0.3cm}
    \begin{alertblock}{统计检验结果}
        \[
        t(8) = -0.024, \quad p = 0.981, \quad d = -0.003
        \]
        \begin{itemize}
            \item $p > 0.05$: 差异在统计学上不显著
            \item $d < 0.2$: 效应量可忽略
        \end{itemize}
    \end{alertblock}
\end{frame}

\begin{frame}{被试级别性能对比}
    \begin{table}[h]
        \centering
        \footnotesize
        \begin{tabular}{lcccc}
            \hline
            \textbf{Subject} & \textbf{Table Q} & \textbf{DQN} & \textbf{$\Delta$} & \textbf{优势方} \\
            \hline
            S01 & 64.47 & 65.71 & -1.24 & DQN \\
            S02 & 58.37 & 56.74 & +1.63 & Table Q \\
            S03 & 68.89 & 71.70 & -2.81 & DQN \\
            S04 & 51.91 & 51.15 & +0.76 & Table Q \\
            S05 & 65.42 & 65.27 & +0.15 & Table Q \\
            S06 & 47.03 & 45.21 & +1.82 & Table Q \\
            S07 & 65.46 & 65.61 & -0.15 & DQN \\
            S08 & 73.56 & 74.62 & -1.06 & DQN \\
            S09 & 68.69 & 68.27 & +0.42 & Table Q \\
            \hline
        \end{tabular}
    \end{table}
    
    \vspace{0.3cm}
    \begin{itemize}
        \item Table Q 在 5/9 被试上优于或持平 DQN
        \item 在低表现被试 (S06) 上优势明显 (+1.82\%)
    \end{itemize}
\end{frame}

\begin{frame}{计算效率对比}
    \begin{table}[h]
        \centering
        \begin{tabular}{lcc}
            \hline
            \textbf{指标} & \textbf{Table Q} & \textbf{DQN} \\
            \hline
            参数量 & 544 (Q 表) & $\sim 10^6$ \\
            训练时间 & \textbf{快 3 倍} & 慢 \\
            推理延迟 & \textbf{低 1000 倍} & 高 \\
            内存占用 & \textbf{低 125 倍} & 高 \\
            收敛性保证 & \checkmark 理论保证 & $\times$ 无保证 \\
            可解释性 & \checkmark 完全可解释 & $\times$ 黑盒 \\
            \hline
        \end{tabular}
    \end{table}
\end{frame}

% ============================================================================
% 消融实验
% ============================================================================
\section{消融实验}

\begin{frame}{Q-Learning 变体对比}
    \begin{table}[h]
        \centering
        \begin{tabular}{lcccc}
            \hline
            \textbf{Method} & \textbf{Acc (\%)} & \textbf{$d$} & \textbf{Kappa} & \textbf{F1} \\
            \hline
            \textbf{Standard Q-Learning} & \textbf{62.68} & 0.00 & 0.502 & 0.620 \\
            Double Q-Learning & 61.86 & -0.10 & 0.491 & 0.610 \\
            Dueling Q-Learning & 62.74 & 0.01 & 0.503 & 0.620 \\
            Dueling Double Q-Learning & 61.73 & -0.11 & 0.489 & 0.609 \\
            DQN & 62.70 & 0.01 & 0.502 & 0.620 \\
            \hline
        \end{tabular}
    \end{table}
    
    \vspace{0.3cm}
    \begin{alertblock}{结论}
        \begin{itemize}
            \item Standard Q-Learning 已是最优选择
            \item Double Q 机制未带来收益 (反而下降 1.3\%)
            \item Dueling 架构效果有限 (+0.1\%)
        \end{itemize}
    \end{alertblock}
\end{frame}

% ============================================================================
% 本章小结
% ============================================================================
\section{本章小结}

\begin{frame}{主要结论}
    \begin{enumerate}
        \item \textbf{自适应策略优于固定策略}: 搜索方法显著优于固定窗口 (+40.6\%)
        \item \textbf{Q-Learning 与基线性能等效}: 与 Grid-Search 差异不显著 ($d < 0.2$)
        \item \textbf{表格型 Q-Learning 可替代 DQN}: $p = 0.981$，性能无显著差异
        \item \textbf{Standard Q-Learning 为最优选择}: 简洁且稳健
        \item \textbf{时间窗选择具有神经生理学一致性}: $t_{\text{start}} \approx 2.2$ s
    \end{enumerate}
\end{frame}
